\section{An explicit subregular reduction of \texorpdfstring{$\mathfrak{sl}_3$}{sl3}}
\label{sec:cbc27-subregular}

The subregular reduction has a nonzero weight-one current and two even charged fields.
Their vacuum products determine both the reduction map and its normalization.
An integral good grading gives a complex with two charged ghost pairs and no neutral ghosts.

Work first over $R=\mathbb C[k]$.
All tensor products and direct sums are algebraic.
The symbol $k$ is an indeterminate of cohomological degree zero.
For matrix units $E_{ij}$, put
\[
 a=E_{12},\quad e=E_{13},\quad A=E_{21},\quad B=E_{31},
 \qquad u=E_{23},\quad v=E_{32},
\]
\[
 h=\operatorname{diag}(1,-1,0),\qquad
 j=\tfrac13\operatorname{diag}(1,1,-2),\qquad
 x=\tfrac13\operatorname{diag}(2,-1,-1),\qquad f=-E_{21}.
\]
Use the invariant form $(s,t)=\operatorname{tr}(st)$ on $\mathfrak{sl}_3$.
The positive eigenspace of $\operatorname{ad}x$ is the abelian algebra
$\mathfrak m=Ra\oplus Re$, in degree one.
The degree-zero space is spanned by $h,j,u,v$, and the degree-minus-one space by $A,B$.
Direct matrix commutation gives
\[
 [x,f]=-f,\qquad
 \mathfrak{sl}_3^f=Rj\oplus Ru\oplus RA\oplus RB.
\]
Thus every eigenvalue on the centralizer is nonpositive.
The character $(f,-)|_{\mathfrak m}$ takes values $-1,0$ on $a,e$.
This is the good-grading convention for quantum reduction in
\cite{SubregularKW027}*{\S1, pp.~3--5}.
Changing the sign of $f$ does not change its nilpotent orbit.

Let $V_R(\mathfrak{sl}_3)$ be the universal affine vacuum vertex algebra with
\[
 [s_\lambda t]=[s,t]+k\operatorname{tr}(st)\lambda\mathbf1.
\]
Tensor it with charged free fermions $b_1,b_2,c_1,c_2$, whose only nonzero brackets are
\[
 [b_i{}_{\lambda}c_j]=[c_j{}_{\lambda}b_i]=\delta_{ij}\mathbf1.
\]
Affine generators have cohomological degree zero, while $|b_i|=-1$ and $|c_i|=1$.
Parity is cohomological degree modulo two.
Write $C_R$ for this differential vertex superalgebra's underlying vertex superalgebra.
Use the mode convention
\[
 Y(s,z)=\sum_{n\in\mathbb Z}s_{(n)}z^{-n-1},\qquad
 [s_\lambda t]=\sum_{n\geq0}\frac{\lambda^n}{n!}s_{(n)}t,
 \qquad :st:=s_{(-1)}t.
\]
All generating modes with $n\geq0$ annihilate the vacuum.
Every state is a finite sum of ordered negative-mode monomials.

Define the odd field and its zero mode by
\begin{equation}\label{eq:cbc27-q}
 q=:(a-\mathbf1)c_1:+:ec_2:,
 \qquad D=q_{(0)}.
\end{equation}
The residue of $z^{-1}dz$ is one.
There are no singular contractions in $q(z)q(w)$, so $D^2=0$.
The zero-mode commutator formula makes $D$ a derivation of every vertex product and gives $D\partial=\partial D$.
In particular,
\begin{align*}
 Db_1&=a-\mathbf1,&Db_2&=e,&Dc_i&=Da=De=0,\\
 Dh&=-2:ac_1:-:ec_2:,&Dj&=-:ec_2:,&Du&=:ec_1:,\\
 Dv&=:ac_2:,&DA&=:hc_1:-:uc_2:+k\partial c_1,\\
 DB&=-:vc_1:+\tfrac12:(h+3j)c_2:+k\partial c_2.
\end{align*}

\subsection{Four closed fields}

Set $N_i=:b_ic_i:$ and define
\begin{equation}\label{eq:cbc27-modified}
 H=h+2N_1+N_2,\qquad J=j+N_2,\qquad
 U=u-:b_2c_1:,\qquad X=v-:b_1c_2:.
\end{equation}
Let $C_R^-$ be the vertex subalgebra generated by $H,J,U,X,A,B,c_1,c_2$.
Its cohomological degrees are nonnegative.

\begin{lemma}\label{lem:cbc27-closed}
The subalgebra $C_R^-$ is preserved by $D$.
Its differential satisfies
\begin{equation}\label{eq:cbc27-doublets}
\begin{gathered}
 DH=-2c_1,\qquad DX=c_2,\qquad DJ=DU=Dc_i=0,\\
 DA=:Hc_1:-:Uc_2:+(k+3)\partial c_1,\\
 DB=-:Xc_1:+\tfrac12:(H+3J)c_2:+(k+3)\partial c_2.
\end{gathered}
\end{equation}
The following four even states are closed:
\begin{equation}\label{eq:cbc27-generators}
 \begin{split}
 J,\qquad U,\qquad
 P&=A+\tfrac14:HH:+:UX:+\tfrac{k+2}{2}\partial H,\\
 V&=B-\tfrac12:HX:-\tfrac32:JX:-(k+2)\partial X.
 \end{split}
\end{equation}
\end{lemma}

\begin{proof}
The affine brackets and the ghost brackets show that $C_R^-$ is a vertex subalgebra with a PBW basis in its displayed generators.
For example,
\[
 [H_\lambda H]=(2k+5)\lambda\mathbf1,\quad
 [H_\lambda J]=\lambda\mathbf1,\quad
 [U_\lambda X]=-\tfrac12H+\tfrac32J+(k+1)\lambda\mathbf1.
\]
The bracket of $A$ or $B$ with a modified current remains a linear combination of $A,B$.
The brackets of modified degree-zero currents close among themselves and scalar terms.
Their brackets with $c_i$ are linear combinations of the $c_i$.

Normal ordering gives the identities
\[
 :N_ic_i:=-\partial c_i,\qquad
 :N_2c_1:+:(:b_2c_1:)c_2:=-\partial c_1,
\]
\[
 :(:b_1c_2:)c_1:+:N_1c_2:=-\partial c_2.
\]
They follow by moving the rightmost ghost mode through the two modes of the normal product.
Substitution in $DA,DB$ proves~\eqref{eq:cbc27-doublets}.
The quasi-commutativity formula gives
\[
 :c_1H:-:Hc_1:=2\partial c_1,
 \qquad :c_1X:-:Xc_1:=-\partial c_2.
\]
Therefore
\begin{align*}
 D\bigl(\tfrac14:HH:\bigr)&=-:Hc_1:-\partial c_1,
 &D(:UX:)&=:Uc_2:,\\
 D\bigl(-\tfrac12:HX:\bigr)&=:c_1X:-\tfrac12:Hc_2:,
 &D\bigl(-\tfrac32:JX:\bigr)&=-\tfrac32:Jc_2:.
\end{align*}
The remaining derivative terms in $DP,DV$ cancel exactly.
\end{proof}

Let $W_R$ denote the vertex subalgebra of $C_R^-$ generated by $J,U,P,V$, with zero differential.

\begin{theorem}[Subregular reduction map]\label{thm:cbc27-subregular-map}
The inclusion
\begin{equation}\label{eq:cbc27-iota}
 \iota:(W_R,0)\lhook\joinrel\longrightarrow(C_R,D)
\end{equation}
is a strict unital quasi-isomorphism of differential vertex superalgebras over $R$.
Ordered normally ordered monomials in all derivatives of $J,U,P,V$ form an $R$-basis of $W_R$.
For every $k_0\in\mathbb C$, specialization gives
\[
 H^r(C_R\otimes_R\mathbb C_{k_0},D)=
 \begin{cases}W_R\otimes_R\mathbb C_{k_0},&r=0,\\0,&r\ne0.\end{cases}
\]
Here $\mathbb C_{k_0}=R/(k-k_0)$, including $k_0=-3$.
The degree-zero vertex algebra is the subregular quantum reduction for $(\mathfrak{sl}_3,x,f)$.
\end{theorem}

\begin{proof}
Let $C_R^+$ be generated by $a-\mathbf1,e,b_1,b_2$.
Its vertex products are commutative and its differential is preserved.
Normally ordered multiplication gives a linear map of complexes
\[
 C_R^-\otimes_R C_R^+\longrightarrow C_R,\qquad
 s\otimes t\longmapsto s_{(-1)}t.
\]
The tensor differential is $D(s\otimes t)=Ds\otimes t+(-1)^{|s|}s\otimes Dt$.
To prove bijectivity, assign filtration degrees zero to derivatives of $a,e,b_i$, one to derivatives of $h,j,u,v,c_i$, and two to derivatives of $A,B$.
All singular brackets lower total filtration degree by at least one.
The associated graded is the polynomial algebra in the even jets tensored with the exterior algebra in the odd jets.
The changes $a\mapsto a-1$ and~\eqref{eq:cbc27-modified}, including all derivatives, are polynomial and invertible there.
Thus multiplication induces a bijection on associated graded modules.
Induction on the nonnegative filtration degree proves the stated linear bijection.
The filtration need not have finite-dimensional pieces.
Every individual state has finite degree.
This is a tensor decomposition of complexes, not of vertex algebras.

Write $a_r'=\partial^r(a-\mathbf1)$, $e_r=\partial^re$, and $b_{ir}=\partial^rb_i$.
The complex $C_R^+$ has differential
\[
 D_+=\sum_{r\geq0}\left(a_r'\frac{\partial}{\partial b_{1r}}+
 e_r\frac{\partial}{\partial b_{2r}}\right).
\]
Odd derivatives act from the left.
For $s_+=\sum_r(b_{1r}\partial_{a_r'}+b_{2r}\partial_{e_r})$, one has
$D_+s_++s_+D_+=N_+$, where $N_+$ counts these factors.
Division by each positive integer eigenvalue gives a contraction onto $R\mathbf1$.
All sums are finite on a state.
Consequently $C_R^-\hookrightarrow C_R$ is a quasi-isomorphism.

On $C_R^-$ assign degree one to derivatives of $H,J,U,X,c_i$, and degree two to derivatives of $A,B$.
The associated graded differential is
\[
 dH=-2c_1,\quad dX=c_2,\quad dJ=dU=dc_i=0,
\]
\[
 dA=Hc_1-Uc_2,\qquad dB=-Xc_1+\tfrac12(H+3J)c_2.
\]
It commutes with all formal derivatives.
Introduce the polynomial variables
\[
 p=A+H^2/4+UX,\qquad
 v_0=B-HX/2-3JX/2.
\]
This triangular change is invertible, and $dp=dv_0=0$.
Hence the graded complex is
\[
 R[\partial^rJ,\partial^rU,\partial^rp,\partial^rv_0:r\geq0]
 \otimes_R
 \bigl(R[H_r,X_r]\otimes\Lambda(c_{1r},c_{2r}),d\bigr).
\]
On the second factor,
\[
 s_-=-\tfrac12\sum_rH_r\frac{\partial}{\partial c_{1r}}
       +\sum_rX_r\frac{\partial}{\partial c_{2r}}
\]
satisfies $ds_-+s_-d=N_-$.
Thus the graded cohomology is precisely the first polynomial factor in degree zero.

The explicit cycles~\eqref{eq:cbc27-generators} lift its four generators.
Their ordered monomials have independent leading symbols, which proves linear independence.
Given a degree-zero closed state, subtract a linear combination of these monomials with the same leading symbol.
The filtration degree decreases.
Repeating terminates and proves that these monomials span $\ker(D|_{(C_R^-)^0})$.
There are no degree-minus-one states in $C_R^-$.
Thus this kernel is $W_R$ itself.
For a closed state in positive cohomological degree, the graded contraction produces a preimage of its leading symbol.
Subtracting its differential lowers the filtration degree and again terminates.
It follows that $H^{r>0}(C_R^-)=0$.

These arguments use only division by nonzero integers and the constant $2$.
They therefore apply over $R$ and after every specialization.
No inverse of $k+3$ enters this proof.
This proves all assertions, including base change and the identification with the defining quantum reduction.
\end{proof}

Theorem~\ref{thm:cbc27-subregular-map} gives an explicit instance of the structure theorem in
\cite{SubregularKW027}*{Theorem~4.1 and its proof, pp.~10--12}.
The displayed contraction proves the required assertions on this universal vacuum carrier directly.
Write $C_k=C_R\otimes_R\mathbb C_k$ and $W_k=W_R\otimes_R\mathbb C_k$ for a fixed complex level.

\subsection{The Bershadsky--Polyakov products}

For $k\ne-3$, set $t=k+3$ and
\begin{equation}\label{eq:cbc27-bp-fields}
 \begin{gathered}
 G^+=U,\qquad G^-=-V,\qquad
 T=\frac{P+\tfrac34:JJ:-\tfrac32\partial J}{t},\\
 \kappa_J=\frac{2k+3}{3},\qquad
 c=-\frac{(2k+3)(3k+1)}{k+3}.
 \end{gathered}
\end{equation}
The sign in $G^-=-V$ fixes the cubic vacuum term.
The fields $G^\pm$ are even.

\begin{proposition}\label{prop:cbc27-bp-products}
The following brackets hold as identities of states in $C_k$, not only in cohomology:
\begin{gather}
 [J_\lambda J]=\kappa_J\lambda\mathbf1,\qquad
 [J_\lambda G^\pm]=\pm G^\pm,\qquad[ G^\pm{}_{\lambda}G^\pm]=0,\label{eq:cbc27-j-products}\\
 [T_\lambda J]=(\partial+\lambda)J,\qquad
 [T_\lambda G^\pm]=(\partial+\tfrac32\lambda)G^\pm,\notag\\
 [T_\lambda T]=(\partial+2\lambda)T+\tfrac{c}{12}\lambda^3\mathbf1.\notag
\end{gather}
\begin{equation}\label{eq:cbc27-g-products}
\begin{split}
 [G^+{}_{\lambda}G^-]
 &=3:JJ:+\tfrac{3(k+1)}2\partial J-tT
   +3(k+1)\lambda J\\
 &\hspace{12pt}+\tfrac{(k+1)(2k+3)}2\lambda^2\mathbf1.
\end{split}
\end{equation}
Consequently $W_k$ is the universal Bershadsky--Polyakov vertex algebra in this normalization.
\end{proposition}

\begin{proof}
Use the noncommutative Wick formula
\[
 [s_\lambda:ab:]=:[s_\lambda a]b:
 +(-1)^{p(s)p(a)}:a[s_\lambda b]:
 +\int_0^\lambda[[s_\lambda a]_\mu b]\,d\mu
\]
and quasi-commutativity $:ab:-(-1)^{p(a)p(b)}:ba:=\int_{-\partial}^0[a_\lambda b]\,d\lambda$.
For $M_{ij}=:b_ic_j:$, their input ghost identities are
\[
 [M_{ij}{}_{\lambda}M_{lm}]
 =\delta_{jl}M_{im}-\delta_{im}M_{lj}
   +\delta_{jl}\delta_{im}\lambda\mathbf1,
\]
\[
 [M_{ij}{}_{\lambda}b_l]=\delta_{jl}b_i,
 \qquad [M_{ij}{}_{\lambda}c_l]=-\delta_{il}c_j.
\]
Affine--ghost mixed brackets vanish.
These formulas immediately give $[J_\lambda J]=\kappa_J\lambda$ and $[J_\lambda U]=U$.
The next substitutions give
\begin{align*}
 [J_\lambda V]&=-V,\\
 [J_\lambda P]&=\tfrac32\lambda J+\tfrac{2k+3}{2}\lambda^2\mathbf1,\\
 [U_\lambda V]&=P-\tfrac94:JJ:-\tfrac{3(k+2)}2\partial J
 -3(k+1)\lambda J-\tfrac{(k+1)(2k+3)}2\lambda^2\mathbf1.
\end{align*}
For the last identity, insert $V=B-:HX:/2-3:JX:/2-(k+2)\partial X$.
The required brackets are
\[
 [U_\lambda B]=A,\quad [U_\lambda H]=U,\quad
 [U_\lambda J]=-U,\quad
 [U_\lambda X]=-H/2+3J/2+(k+1)\lambda.
\]
The two Wick integrals combine to
$\lambda U_{(0)}X+\frac{k+1}{2}\lambda^2\mathbf1$.
Expanding the four displayed summands, and using quasi-commutativity for $H,J$, yields the asserted formula.
Substitution of~\eqref{eq:cbc27-bp-fields} proves~\eqref{eq:cbc27-g-products}.
The formula for $[J_\lambda P]$ and $[J_\lambda:JJ:]=2\kappa_J\lambda J$ give the stated $T,J$ bracket.

For the conformal normalization, let $S$ be the affine Sugawara field and put
\[
 L^{\mathrm{full}}=S+\partial x+\sum_{i=1}^2:(\partial b_i)c_i:.
\]
Its affine generator weights are $1-\deg_x(s)$.
Each ghost pair has weights $(0,1)$ and central charge $-2$.
The field $q$ is primary of weight one, so $DL^{\mathrm{full}}=0$.
The Sugawara and ghost contractions give
\[
 c_L=\frac{8k}{t}-8k-4,
 \qquad [L^{\mathrm{full}}{}_{\lambda}J]
 = (\partial+\lambda)J-\frac{\kappa_J}{2}\lambda^2\mathbf1.
\]
Here $(x,x)=2/3$, $(x,j)=1/3$, and the second ghost pair contributes $-\lambda^2/2$.
There is an explicit homotopy identity:
\begin{equation}\label{eq:cbc27-stress-homotopy}
 t\bigl(L^{\mathrm{full}}-T-\tfrac12\partial J\bigr)
 =D\bigl(:b_1A:+:b_2B:\bigr).
\end{equation}
For its verification, write the affine Sugawara numerator as
\[
 tS=\tfrac14:hh:+\tfrac34:jj:
 +\tfrac12\bigl(:aA:+:Aa:+:eB:+:Be:+:uv:+:vu:\bigr).
\]
Use
$:Aa:=:aA:-\partial h$,
$:Be:=:eB:-\partial(h+3j)/2$, and
$:vu:=:uv:+\partial(h-3j)/2$.
Substitute~\eqref{eq:cbc27-modified}--\eqref{eq:cbc27-generators}.
The remaining terms are
$:(a-1)A:+:eB:-:b_1DA:-:b_2DB:$,
which equal the right side of~\eqref{eq:cbc27-stress-homotopy} by the odd derivation rule.

Improving $L^{\mathrm{full}}$ by $-\partial J/2$ changes the central charge by $3\kappa_J$.
It makes $J$ primary and gives $G^\pm$ weights $3/2$.
Here charge is the eigenvalue of $J_{(0)}$.
Under $L^{\mathrm{full}}_{(1)}$, the states $H,J,U,X$ have weight one and $A,B,P,V$ have weight two.
The improvement adds half the charge to these weights.
Its cohomology class is $[T]$.
The inclusion $W_k\to C_k$ is injective on cohomology.
Thus its conformal identities hold strictly in $W_k$.
The value $c=c_L+3\kappa_J$ is exactly~\eqref{eq:cbc27-bp-fields}.
Higher products $T_{(n)}G^\pm$, $n\geq2$, have charge $\pm1$ and weight at most $1/2$.
The PBW basis has no such states.
Products $G^\pm{}_{(n)}G^\pm$, $n\geq0$, have charge $\pm2$ and weight at most two.
The smallest PBW weight with either charge is three.
These products therefore vanish.
This proves all brackets.
They agree with the universal presentation in
\cite{SubregularArakawa027}*{\S2, p.~2}.
The PBW theorem proves that this presentation introduces no further relations.
\end{proof}

\begin{corollary}[Special levels and reflection]\label{cor:cbc27-bp-levels}
The current vacuum coefficient is $\kappa_J=(2k+3)/3$.
The $G^+,G^-$ cubic coefficient is $(k+1)(2k+3)$, and the $T,T$ quartic coefficient is $c/2$.
All three vanish simultaneously at $k=-3/2$.
At $k=-1/2$ they are $2/3,1,1/5$, respectively.
For $k'=-k-6$ and $k\ne-3$, the reflected conformal charges obey
\[
 c(k)+c(k')=50.
\]
Equality $c(k)=c(k')$ holds precisely when $(k+3)^2=-4$.
Hence a conformal vertex isomorphism at reflected levels is impossible for all other noncritical parameters.
\end{corollary}
\begin{proof}
The common root follows from the first coefficient and direct substitution in the other two.
At $k=-1/2$, the charge is $2/5$.
Writing $t=k+3$ gives $c(k)=25-6t-24/t$ and $c(k')=25+6t+24/t$.
Their equality is equivalent to $t^2=-4$.
A conformal unital map preserves $T_{(3)}T=(c/2)\mathbf1$.
At the two exceptional parameters, this calculation states only equality of charges.
It does not construct an isomorphism.
\end{proof}

\subsection{Unital binary collisions}

Fix a complex level $k$, including $k=-3$, and an affine coordinate $z$ on $\mathbb A^1$.
The map $\iota:W_k\to C_k$ commutes with translation and every vertex product.
For a differential vertex algebra $F$, its right differential-operator module is
$\mathcal A_F=\omega_{\mathbb A^1}\otimes_{\mathbb C}F$, with
\[
 (g(z)\,dz\otimes a)\cdot\partial_z
 =-g'(z)\,dz\otimes a-g(z)\,dz\otimes\partial a.
\]
In left-module notation, put
\[
 \mathcal P_F=\mathbb C[z,w,(z-w)^{-1}]\otimes F\otimes F,
 \qquad
 \mathcal R_F=F\otimes\mathbb C[w]\otimes \xi^{-1}\mathbb C[\xi^{-1}],
 \quad \xi=z-w.
\]
The source connections are $\partial_z^{\mathrm{base}}+\partial_F\otimes1$ and
$\partial_w^{\mathrm{base}}+1\otimes\partial_F$.
The target actions are
\[
 \partial_z=\partial_\xi,\qquad
 \partial_w=\partial_w^{\mathrm{base}}-\partial_\xi+\partial_F.
\]
Multiplication by $\xi$ on the target is taken modulo regular powers.
Tensoring with $dz\wedge dw$ gives the corresponding right modules.
The orientation is $dz\wedge dw=d\xi\wedge dw$, with $\operatorname{Res}\xi^{-1}d\xi=1$.
The tensor differential on $\mathcal P_F$ uses the sign $(-1)^{|a|}$ on its second state.
Define the full meromorphic collision by
\[
 \mu_F\bigl(g(z,w)\otimes a\otimes b\bigr)
 =\operatorname{pp}_{\xi=0}\bigl(g(w+\xi,w)Y(a,\xi)b\bigr).
\]
Here $\operatorname{pp}$ retains the strictly negative powers of $\xi$.
The answer is finite by vertex truncation and the finite pole order of $g$.
Both $\mu_F$ and $\iota$ commute with the internal differential.

\begin{proposition}\label{prop:cbc27-binary}
The map $\mathcal A_{W_k}\to\mathcal A_{C_k}$, $g\,dz\otimes a\mapsto g\,dz\otimes\iota(a)$,
is a unital quasi-isomorphism of differential chiral algebras on this affine chart.
Its tensorwise maps make the square of $\mu_{W_k}$ and $\mu_{C_k}$ commute exactly.
They are quasi-isomorphisms on $\mathcal P$ and $\mathcal R$, and on their collision mapping cones.
In particular, the vacuum terms are preserved:
\[
 \mu_{C_k}(J\otimes J)=\kappa_J\xi^{-2}\mathbf1,
\]
and, for $k\ne-3$,
\[
 \begin{split}
 \mu_{C_k}(G^+\otimes G^-)
 ={}&(k+1)(2k+3)\xi^{-3}\mathbf1
 +3(k+1)\xi^{-2}J\\
 &+\xi^{-1}\bigl(3:JJ:+\tfrac{3(k+1)}2\partial J-(k+3)T\bigr).
 \end{split}
\]
\end{proposition}
\begin{proof}
Equality of every mode product gives the commutative square coefficient by coefficient.
The identities $Y(\partial a,\xi)=\partial_\xi Y(a,\xi)$ and
$[\partial_F,Y(a,\xi)]=\partial_\xi Y(a,\xi)$ prove compatibility with both connections.
The vertex skew-symmetry and Jacobi identities give the chiral identities after taking principal parts and converting to right modules.
The vacuum embedding is $g\,dz\mapsto g\,dz\otimes\mathbf1$.
Its collision is $\operatorname{pp}_{\xi}(\xi^{-1}Y(\mathbf1,\xi)a)=\xi^{-1}a$.
Localization and tensor products over $\mathbb C$ preserve quasi-isomorphisms.
The same holds for the direct sum $\xi^{-1}\mathbb C[\xi^{-1}]$.
For a cochain map $\mu:P\to R$, use
$\operatorname{Cone}(\mu)^n=R^n\oplus P^{n+1}$ and
$d(r,p)=(d_Rr+\mu p,-d_Pp)$.
The cone's short exact sequence gives the cone quasi-isomorphism.
The two displayed collisions follow directly from~\eqref{eq:cbc27-j-products}--\eqref{eq:cbc27-g-products}.
The same coefficient argument preserves every finite iterated vertex product.
\end{proof}

There is no strict unital vertex map $C_k\to W_k$ preserving cohomological degree.
It would kill $b_i,c_i$ and contradict $b_i{}_{(0)}c_i=\mathbf1$.
The inverse to $\iota$ is therefore an inverse in the category obtained by inverting quasi-isomorphisms.
The polynomial family and its binary collision maps specify no comparison with the reflected Koszul dual.
A completed chiral bar, a second differential, and descent through nodal curves require their own constructions.
